\chapter{Introduction} \label{sec:intro}
The aim of this guide is to give to future FPGA software developers some basic knowledge to deal with the installation of Petalinux SDK and the creation of the environment needed to develop ad-hoc root file systems for application concerning SoM(s) and more. As a matter of fact, this guide has been written with the goal to give a general procedure to install different version of Petalinux on different (compatible) OS. All the files used are to be referred to the AMD website. In particular, all the needed libraries can be found in the file at this \href{https://adaptivesupport.amd.com/s/article/000034483?language=en_US}{page}. A more efficient and quick way to install them in given in the third section.

However, since the installation and build processes are strictly related to the project and to the environment, it is impossible for this guide to collect and solve all the errors that might be encountered. It is thus recommended that the the person approaching this field possesses sufficient problem-solving skills so as to have the ability to deal with the bespoken errors by oneself. A warning that has to be given is that the use of Petalinux and WSL is a littla less trivial than using directly a compatible OS. If one chooses to follow this approach, not-well-documented bugs my cause the failure of of the processes listed in the following sections. 